\section{Dominio positivo de la coordenatización elíptica}
\begin{proposition}[Pertenencia de las representaciones a la coordenatización elíptica]
\label{iv:ellipse:pertenencia}
La rama positiva de \eqref{eq:ii-par-angular}, con los
intervalos \eqref{md:accion:intervalos}, satisface
\(A>0\), \(0<C^*<A\) y \(\hbar_s>0\).
La reversión de la coordenada impar conserva \(|C^*/A|<1\).
\end{proposition}
\begin{proof}
La proposición \ref{md:accion:positividad} da
\(0<\eta_{\mathrm{ret}}<H_5\) y la positividad de ambas
secciones de acción. Para una cota superior de \(H_5\),
\eqref{md:accion:intervalos} implica
\(\alpha<1/125\) y \(\varphi>8/5\). Al conservar sólo
los términos positivos de \eqref{eq:el-h5}, se obtiene
\[
 H_5<
 \frac{\sqrt5}{2}+
 \frac{9}{16\cdot125^2}+
 \frac{7}{48\cdot125^4}.
\]
Como \(5<81/16\), se tiene \(\sqrt5/2<9/8\). Además,
\[
 \frac{9}{16\cdot125^2}+
 \frac{7}{48\cdot125^4}
 =\frac{421882}{11718750000}<\frac1{1000}.
\]
Por tanto, \(H_5<9/8+1/1000<2\). Usando ahora
\(\alpha>7/1000\),
\[
 0<C^*=2(\eta_{\mathrm{ret}}+\alpha)
 <2\left(2+\frac1{125}\right)=\frac{502}{125}
 <7<1000\alpha=A.
\]
El intercambio de canales de
\ref{prop:ii-inversion-canales-angulares} cambia \(C^*\) por \(-C^*\)
y conserva \(A\); por ello conserva la desigualdad absoluta.
\end{proof}
